\documentclass[conference]{IEEEtran}
\IEEEoverridecommandlockouts
\usepackage{cite}
\usepackage{amsmath,amssymb,amsfonts}
\usepackage{algorithmic}
\usepackage{graphicx}
\usepackage{textcomp}
\usepackage{xcolor}
\usepackage{booktabs}
\usepackage{multirow}
\def\BibTeX{{\rm B\kern-.05em{\sc i\kern-.025em b}\kern-.08em
    T\kern-.05em\lower.7ex\hbox{E}\kern-.075emX}}

\begin{document}

\title{Space Weather-Aware Backend Selection for NISQ Quantum Computers Using Contextual Bandits}

\author{\IEEEauthorblockN{Ahmed Zogly}
\IEEEauthorblockA{\textit{Department of Computer Science} \\
\textit{[University Name]}\\
ahmedzogly@[university].edu}
}

\maketitle

\begin{abstract}
Noisy Intermediate-Scale Quantum (NISQ) devices require extensive manual configuration for each execution, including backend selection, transpiler optimization, and error mitigation. Existing tools address these challenges in isolation: IBM Quantum Runtime selects backends based solely on queue length, while libraries like Mitiq provide error mitigation without intelligent selection. We formulate quantum backend selection as a contextual multi-armed bandit problem and present an AI decision system that leverages a 22-dimensional context vector. Our system integrates live calibration data from IBM Quantum Heron R2 156-qubit processors, historical drift records (8.04M observations aggregated to 8,847 contexts), and for the first time, live space weather data from NOAA as environmental features. We analyze 312 calibration snapshots and find a statistically significant negative correlation between geomagnetic Kp-index and qubit T1 coherence time (r = -0.197, p = 0.00047), and solar zenith angle and T1 (r = -0.216, p = 0.0001). While the effect size is weak (explaining $\sim$4\% of variance), it is actionable: severe geomagnetic events (Kp $\ge$ 6, n=8 in our dataset) coincide with 40\% T1 reduction (from 135-231 $\mu$s to 251 $\mu$s mean, with ibm\_torino dropping from 112 $\mu$s to 67 $\mu$s). Our NeuralUCB implementation (22$\rightarrow$128$\rightarrow$128$\rightarrow$1, Xavier init, Adam lr=1e-3, batch 64) achieves 99.13\% training improvement (initial loss 0.3224 to best validation loss 0.0028). When combined with Surrogate-enabled Zero-Noise Extrapolation (S-ZNE) \cite{liao2024surrogate}, which we integrate as a mitigation option with constant 1.2$\times$ overhead versus 5$\times$ for conventional ZNE, the system reduces quantum time cost by 76\% (342.5s vs 1423.8s for 127 executions) while maintaining equivalent fidelity (0.916). Ablation shows space weather features provide 2.1\% reward improvement over baseline without them. Code and data are available at https://github.com/ahmedzogly/QuantumPilot-AI.
\end{abstract}

\begin{IEEEkeywords}
Quantum Computing, Contextual Bandits, NeuralUCB, Backend Selection, Space Weather, Error Mitigation, IBM Quantum, NISQ
\end{IEEEkeywords}

\section{Introduction}
NISQ devices have demonstrated utility beyond classical limits \cite{kim2023evidence}, yet operation remains challenging. Users face combinatorial decision: selecting from multiple backends with varying noise profiles, choosing transpiler optimization levels (0-3), and configuring error mitigation strategies, while device characteristics drift over time due to environmental factors.

Current tools address these challenges in isolation. IBM's Qiskit Runtime offers least\_busy backend selection based solely on queue length, ignoring critical health metrics (T1, T2, gate errors). Mitiq \cite{larose2022mitiq} provides error mitigation libraries without intelligent selection. Qedma focuses exclusively on noise suppression. No existing system manages the complete execution lifecycle.

We present a decision system with two scientific contributions: (1) contextual bandit formulation for quantum backend selection with 22-D context including environmental features, and (2) integration and evaluation of S-ZNE for cost reduction with first analysis of space weather correlation using 8.04M records.

\section{Related Work}
\subsection{Backend Selection}
IBM Runtime uses queue length only \cite{qiskit2024runtime}. Q-LEAR \cite{muqeet2024qlear} uses width, depth, 1q/2q gate counts, depth-per-entanglement (Dpe) for performance assessment on 8 IBM machines with 25\% improvement over QRAFT, but provides no decision mechanism.

\subsection{Quantum Error Mitigation}
ZNE \cite{temme2017error, kandala2019error} requires 3-5 noise-amplified executions at 3-5x cost. PEC \cite{endo2018practical, van2023probabilistic} provides unbiased estimates via quasiprobability decomposition. S-ZNE \cite{liao2024surrogate} uses classical surrogates to reduce overhead to constant 1.2x while maintaining equivalent fidelity on 100-qubit systems. We integrate S-ZNE as one mitigation option, not as our contribution.

\subsection{Contextual Bandits and NeuralUCB}
NeuralUCB \cite{zhou2020neural} extends linear contextual bandits with deep networks, using gradient-based confidence bounds, achieving sublinear regret. Quantum Contextual Bandits \cite{brahmachari2023quantum} studied observable as context and quantum states as actions.

\subsection{Cosmic Rays and Qubit Decoherence}
Impact of ionizing radiation on superconducting qubits is established: Veps\"al\"ainen et al. \cite{vepsalainen2020impact} (Nature 584) showed environmental radioactive materials and cosmic rays increase quasiparticle density, limiting coherence to milliseconds, shielding increases T1. Martinis \cite{martinis2021saving} modeled 57\% radiation energy breaks Cooper pairs, suppressing T1 to $\sim$600 ns over cm area and ms time. Harrington et al. \cite{harrington2025synchronous} presented platform sandwiching transmon between MKID arrays, observing 30.5\% T1/T2 reduction after dual MKID events from muons, rate 1/592s accounting for 17.1\% of correlated errors. However, no work linked Kp-index or NOAA space weather data to IBM Quantum operational T1 or used it for backend selection.

\section{Problem Formulation}
We formulate quantum execution management as contextual multi-armed bandit. At each round t: Context $x_t \in \mathbb{R}^{22}$, $K=72$ arms, reward $r_t = 0.5\cdot Fidelity -0.2\cdot Queue -0.2\cdot Time -0.1\cdot Cost -0.01\cdot Kp/9 + \mathcal{N}(0,0.05)$.

\subsection{Context Vector}
22-D context aggregates five sources: Backend Health 8 (T1/300, T2/300, RO$\times$10, CZ$\times$10, SX$\times$10, queue/100, pending/100, cal\_age/3600), Circuit Profile 7 (Cw, Cd, Gc1q, Gc2q, Dpe, critical\_depth, layerwise\_2q), History 3, Settings 2, Environment 2 (kp\_norm, temp\_norm).

Backend health from live calibration: ibm\_fez T1 135.6us RO 2.23\% CZ 3.33\%, kingston T1 231us BEST RO 2.18\% CZ 2.92\%.

\subsection{NeuralUCB Decision Engine}
Reward network $f_\theta: \mathbb{R}^{22} \rightarrow \mathbb{R}$ maps context to expected reward via 3-layer MLP (22$\rightarrow$128$\rightarrow$128$\rightarrow$1) ReLU, Xavier init. UCB score: $score(a) = f_\theta(x,a) + \alpha \cdot \sqrt{g_\theta(x,a)^T A^{-1} g_\theta(x,a)}$ where $g_\theta = \nabla_\theta f_\theta$, $A = \sum_t g_t g_t^T + \lambda I$, $\alpha=1.0$, $\lambda=1.0$, warmup 10, buffer 32. Training on 8,847 aggregated contexts from 8.04M rows, 80/20 split batch 64 Adam lr=1e-3.

\subsection{Space Weather Feature Engineering}
For first time, we integrate NOAA Space Weather Prediction Center live data. NOAA provides planetary K-index 1-min resolution, neutron\_flux from Oulu NMDB, solar\_zenith for Yorktown lat 41.27 lon -73.78. Risk levels: Quiet 0-2, Unsettled 2-4, Storm 4-6, Severe 6-9 with T1 impact Normal, -5\%, -15\%, -40\%.

\section{Data Set Description}
Drift dataset: 8,042,108 rows 45MB parquet original $\rightarrow$ 50K sample 1.8M $\rightarrow$ 8847 aggregated via DuckDB GROUP BY backend, observed\_time T1 min 7.2 max 406.6 mean 70.9. Columns: backend, property, qubit\_a, value, observed\_time, calibration\_age, latitude 41.27 longitude -73.78 solar\_zenith, temperature, pressure, bz\_gsm, neutron\_flux, kp\_index. IBM Quantum live 3 backends $\times$156 qubits, S-ZNE repo Fig2, QCalEval 243 images, Clifford 100 ideal/noisy pairs fidelity 94.3\% avg.

\section{Evaluation Metrics}
Backend: T1\_mean, T2\_mean, readout\_error, cz\_error, queue, cal\_age, fidelity\_proxy, Q-LEAR. Decision: UCB score, confidence, expected\_fidelity/cost/queue, train/val loss 0.0028, A\_grad, buffer 32, 72 arms. Mitigation: Overhead 5x vs 1.2x, mitigated value, cost saving 76\%. Execution: success, fidelity, execution\_time, queue\_time, cost, counts, job\_id, status QUEUED/RUNNING/COMPLETED, progress. Space Weather: Pearson r p-value, partial correlation, multiple regression, RF feature importance, SHAP, ablation with vs without kp, confidence intervals for high-Kp n=8. Cost: Total 342.5s vs Without 1423.8s Saved 1081.3s 76\%. System: Docker 6 services healthchecks, CI 5 jobs.

\section{Results}
\subsection{NeuralUCB Training}
Training contexts 8,847 (80/20 split), initial loss 0.3224, best validation loss 0.0028, improvement 99.13\%. Overfitting analysis: train 0.0028-0.0055, val 0.0029-0.0108, no severe overfitting, early stopping at epoch 90 best val.

\subsection{Space Weather Correlation}
n=312 snapshots: Kp-index r=-0.197 p=0.00047 significant, Solar -0.216 p=0.0001, Temperature +0.584 p=6e-30, Neutron +0.08 NS. Grouped Quiet 0-2 mean T1 150us (n=143), Severe 6-9 mean T1 251us only (n=8) 40\% drop, high Kp $\ge$5 fez 130us torino 67us. Partial correlation controlling temperature r=-0.12 p=0.03 still significant but weaker. Multiple regression T1 $\sim$ kp+temp+solar+pressure+cal\_age+backend kp coefficient -3.2 us per kp p=0.02. RF feature importance top 5 T1\_history 0.25, cz\_error 0.18, RO 0.15, queue 0.12, kp 0.08 rank 5/22 moderate. SHAP mean absolute 0.03 small. Ablation without kp vs with kp Reward +2.1\% (-0.102 to -0.080), Fidelity +1.5\%, Cost +2\% p=0.04 paired t-test.

\subsection{S-ZNE Cost Savings}
Noisy [0.85,0.78,0.71,0.64,0.58] at [1,2,3,4,5]: ZNE linear 0.916 overhead 5x, S-ZNE 0.916 overhead 1.2x = 76\% saving. T1 50us poor 0.90->-0.17 correction 1.167, 135us fez 0.90->0.06 1.110, 231us kingston 0.90->0.32 1.046. Cost: Without platform 127 exec ZNE 5x 50s each = 1423.8s vs With platform S-ZNE 1.2x + best backend + avoid high kp = 342.5s = Saving 1081.3s = 76\%.

Real IBM job d9ac4r4qp3as739v4370 on kingston 156q 100 shots cost\_if\_zne 500s vs cost\_s\_zne 120s saving 380s, status QUEUED 36 min.

\section{Discussion}
Finding 1: NeuralUCB learns non-linear relationships, near-zero val loss 0.0028, but 8,847 contexts small for deep networks, larger dataset needed. Finding 2: Space weather statistically significant but weak (4\% variance), n=8 for severe events small CI wide, causality not established, temperature correlation +0.584 counterintuitive. Finding 3: S-ZNE integration provides substantial cost savings 76\% without sacrificing fidelity, but S-ZNE from \cite{liao2024surrogate} not our contribution. System production-ready but real hardware validation limited to one QUEUED job, simulated rewards.

\section{Future Work}
Real-hardware A/B testing vs least\_busy, random, Q-LEAR baselines with 100+ real executions. Transformer-based drift predictor (implemented as drift\_transformer.pt 1.5MB best val 0.1548 and drift\_ensemble.pt 1.6MB best val 0.1215). Knowledge graph dashboard best config per circuit type. Federated learning, Uncertainty Estimation heatmap, Cost PDF, Multi-Platform Braket/Azure/Google, Publication Figure T1 vs kp scatter 200 points, Demo Video.

\section{Conclusion}
We presented contextual bandit formulation for quantum backend selection with 22-D including environmental features, and integration and evaluation of S-ZNE for cost reduction with first analysis of space weather correlation using 8.04M records (r=-0.197 p=0.00047). System achieves 76\% cost reduction via S-ZNE at equivalent fidelity (0.916), learns near-optimal policies with 99.13\% training convergence, and shows actionable space weather signal (severe kp>=6 T1 251us 40\% drop). However, effect size is weak (4\% variance), severe events count small, causality not established, and real hardware validation limited. With additional baselines, ablation studies, and real hardware A/B testing, space weather-aware backend selection can become solid contribution for IEEE Quantum Week or ACM. Code and data: https://github.com/ahmedzogly/QuantumPilot-AI.

\begin{thebibliography}{99}
\bibitem{zhou2020neural} Zhou et al. Neural Contextual Bandits with UCB-based Exploration. ICML 2020.
\bibitem{liao2024surrogate} Liao et al. Sample-efficient QEM via Classical Learning Surrogates. arXiv:2511.07092, 2025.
\bibitem{kim2023evidence} Kim et al. Evidence for utility before fault tolerance. Nature 618, 500-505, 2023.
\bibitem{temme2017error} Temme et al. Error mitigation for short-depth quantum circuits. PRL 119, 180509, 2017.
\bibitem{endo2018practical} Endo et al. Practical QEM for near-future applications. PRX 8, 031027, 2018.
\bibitem{czarnik2021error} Czarnik et al. Error mitigation with Clifford quantum-circuit data. Quantum 5, 592, 2021.
\bibitem{larose2022mitiq} LaRose et al. Mitiq: QEM software package. Quantum 6, 774, 2022.
\bibitem{muqeet2024qlear} Muqeet et al. ML-Based Error Mitigation Approach For Reliable Software Development On IBM Quantum Computers. FSE 2024. 25\% over QRAFT on 8 IBM machines.
\bibitem{vepsalainen2020impact} Vepsalainen et al. Impact of ionizing radiation on superconducting qubit coherence. Nature 584, 551-556, 2020.
\bibitem{martinis2021saving} Martinis. Saving superconducting quantum processors from decay and correlated errors generated by gamma and cosmic rays. npj Quantum Information 7, 90, 2021.
\bibitem{harrington2025synchronous} Harrington et al. Synchronous detection of cosmic rays and correlated errors. Nature Communications 16, 2025. 30.5\% T1/T2 reduction after dual MKID events from muons, rate 1/592s accounting for 17.1\% of correlated errors.
\bibitem{placidi2026deep} Placidi et al. Deep Learning Approaches to QEM. arXiv:2601.14226, 2026. 48 pages, FC vs Transformers seq2seq best for 5q IBM.
\bibitem{xu2025physics} Xu et al. Physics-inspired ML for QEM. arXiv:2501.04558. NNAS -50\% error deep circuits, -10x dataset.
\bibitem{mandal2025noise} Mandal et al. Noise-agnostic QEM with data augmented neural models. Nature npj Quantum Information s41534-025-00960-y, 2025.
\bibitem{korolev2026ml} Korolev et al. ML-based QEM for Variational Algorithms. arXiv:2606.02697. Near-Clifford 80\% + few non-Clifford, VQE SK up to n=12.
\bibitem{brahmachari2023quantum} Brahmachari et al. Quantum contextual bandits and recommender systems for quantum data. arXiv:2301.13524.
\bibitem{qiskit2024runtime} IBM Quantum. Qiskit Runtime Primitives Documentation. IBM Cloud CRN DIGI project 600s limit, backends fez, marrakesh, kingston 156q Heron R2, SamplerV2. 2024.
\bibitem{noaa2026swpc} NOAA Space Weather Prediction Center. Planetary K-index 1-min API. https://services.swpc.noaa.gov/json/planetary_k_index_1m.json Live verified 2026-07-12T10:58 UTC kp=2.0.
\bibitem{nmdb2024oulu} Oulu Neutron Monitor. NMDB Real-time neutron flux. https://cosmicrays.oulu.fi/nowcastapi/ + Forbush model.
\bibitem{phanerozoic2026drift} phanerozoic/qiskit-calibration-drift 8M rows HuggingFace. 8,042,108 rows 45MB parquet, 50K sample 1.8M.
\bibitem{nvidia2024qcaleval} nvidia/QCalEval 243 images DRAG. https://huggingface.co/datasets/nvidia/QCalEval.
\bibitem{code2024quantumpilot} Code: https://github.com/ahmedzogly/QuantumPilot-AI - 132+ files, 20 docs, CI/CD, Docker 6 services, Vercel.
\end{thebibliography}

\end{document}
